% Optional companion theorem. All labels use ert:.
\section{Effective eventual reconstruction for every fixed even order}
\label{ert:section}

The structural factorization used for order eight has an all-order
consequence. It reduces the full conjecture at each fixed smaller
order to an explicitly bounded finite calculation. It does not prove
that all those finitely many matrices are nonsingular.

\begin{theorem}[Effective eventual optimal-prefix rank]
\label{ert:main}
For every integer $a\ge1$ there is a constructible integer $B_a$
such that the optimal-prefix coefficient matrix $M_{a,b}$ is
nonsingular for every integer $b\ge B_a$. Thus, for every fixed
even tail order $p=2a$, only finitely many larger even orders can
possibly fail the source optimal-prefix conjecture. The formulas
below give a terminating algorithm for an integer $B_a$, using
finite polynomial arithmetic and exact integer comparisons.
\end{theorem}

Here and below, the general source matrix is
\begin{equation}\label{ert:M}
 (M_{a,b})_{ik}=B_{2(i+k-a+1)}
 \left\{\frac{\binom{2i+2}{2a-2k}}{2i+2}
 +\frac{\binom{2b+2i-2a+2}{2b-2k}}{2b+2i-2a+2}\right\},
 \quad0\le i,k<b.
\end{equation}
Negative Bernoulli indices give zero; each binomial summand is
omitted unless its lower index lies between $2$ and its upper index.
This is the same normalization as \eqref{r8c:M} at $a=4$.

\subsection{Exact normalization of the retained determinant}

Assume initially $b\ge2a$. For $0\le i,\ell<a$ define
\[
 n_i=b-a+i,\quad N_{i\ell}=2(n_i-\ell)+1,\quad
 R_\ell=2b-2a+2+2\ell,
\]
\[
 F_{i\ell}=\frac1{(2\ell+1)!}
 \prod_{r=0}^{2\ell-1}(2n_i-r)(2n_i+3+r),
\]
and define integer-polynomial coefficients $d_{ih}$ by
\begin{equation}\label{ert:newton}
 (2n_i+2-y)(4n_i+2-y)^{\underline{2i}}
       =\sum_{h=0}^{2i+1}d_{ih}y^{\underline h}.
\end{equation}
They lie in $\mathbb Z[b]$: the change of basis from ordinary
powers to falling powers has integer Stirling-number coefficients.
Let $U_1,\ldots,U_{2a-1}$ be independent variables and put
\begin{equation}\label{ert:A}
 \mathcal A_{i\ell}=2F_{i\ell}+
 \frac{R_\ell!}
 {(2n_i+1)(2n_i+2)(2\ell+2)!N_{i\ell}!}
 \sum_{k=0}^{i}d_{i,2k+1}N_{i\ell}^{\underline{2k+1}}
            U_{a-i+\ell+k}.
\end{equation}
The inherited reduction is the theorem ``Reduction to a determinant
of size $a$,'' equation \texttt{eq:general-small-det}, in the source
file \texttt{03b\_general\_hankel.tex}~\cite{GeneralHankel}.
We give its exact multiplier and the normalization argument, so
that no nonvanishing statement about an uncomputed inverse matrix
is needed.

\begin{lemma}[Exact source-to-retained determinant multiplier]
\label{ert:normalization}
Put
\[
 h_n=\frac{((2n+1)!)^2}{(4n+2)!(4n+3)!},\qquad
 c_n=\frac{(2n+1)(2n+2)!}{(4n+2)!},
 \qquad \epsilon_{a,b}=a(b-a)+\frac{a(a-1)}2.
\]
For $b\ge2a$, with $U_j=2B_{2b-2j}$,
\begin{equation}\label{ert:exact-multiplier}
 \det M_{a,b}=(-1)^{\epsilon_{a,b}}
 \frac{\prod_{t=0}^{b-1}(2b+2t-2a+1)!}
      {\prod_{k=0}^{b-1}(2b-2k)!}
 \left(\prod_{n=0}^{b-a-1}h_n\right)
 \left(\prod_{i=0}^{a-1}c_{n_i}\right)\det\mathcal A.
\end{equation}
Every scalar factor other than the displayed sign is strictly
positive. In particular,
\begin{equation}\label{ert:equivalence}
 \det M_{a,b}\ne0
 \quad\Longleftrightarrow\quad
 \det\mathcal A\bigl(b;2B_{2b-2},\ldots,2B_{2b-4a+2}\bigr)\ne0.
\end{equation}
\end{lemma}

\begin{proof}
Set $\mu_j=B_{2j}/(2j)!$ for $j\ge0$ and $\mu_j=0$ for $j<0$.
Cancelling factorials in \eqref{ert:M} gives
\[
 M_{a,b}=R(H+U^{\mathrm{pert}})E,
 \quad R_{tt}=(2b+2t-2a+1)!,\quad E_{kk}=\frac1{(2b-2k)!},
\]
where $H_{tk}=\mu_{t+k-a+1}$ and
\[
 U^{\mathrm{pert}}_{tk}=\mathbf1_{k<a}\mu_{t+k-a+1}
 \frac{(2t+1)!(2b-2k)!}
 {(2a-2k)!(2b+2t-2a+1)!}.
\]
The diagonal symbols $R,E$ in this proof are matrices; $R_\ell$
in \eqref{ert:A} is the previously defined integer.

The explicit monic Pad\'e rows are
\[
 \pi_n(x)=\sum_{t=0}^np_{nt}x^t,\qquad
 p_{nt}=\frac{2(2n+2t+1)!(2n+1)!}
 {(4n+2)!(2n-2t+1)!(2t)!}.
\]
For $\mathcal N(x^r)=\mu_{r+1}$ they satisfy
\begin{equation}\label{ert:pade-orthogonality}
 \mathcal N(x^j\pi_n)=0\ (0\le j<n),\qquad
 \mathcal N(x^n\pi_n)=h_n.
\end{equation}
Here is a finite verification of the normalization. Define
\[
 P_N(t)=\sum_{k=0}^N
 \frac{(2N-k)!N!}{(2N)!k!(N-k)!}t^k.
\]
Direct coefficient comparison gives
\[
 e^tP_N(-t)-P_N(t)
 =\frac{(-1)^N(N!)^2}{(2N)!(2N+1)!}t^{2N+1}
       +O(t^{2N+2}).
\]
Indeed, the coefficients through $2N$ reduce to the finite identity
$\sum_k(-1)^k\binom Nk(1+x)^{2N-k}=x^N(1+x)^N$;
the next coefficient is
$(2N)!^{-1}\int_0^1x^N(x-1)^N\,dx$.
Taking $N=2n+1$, set
\[
 q_n(t^2)=\frac{P_N(t)-P_N(-t)}t,\qquad
 r_n(t^2)=\frac{P_N(t)+P_N(-t)}2.
\]
For $\Phi(z)=\sum_{j\ge0}\mu_jz^j
=\frac{\sqrt z}{2}\coth\frac{\sqrt z}{2}$ this proves
\[
 \Phi(z)q_n(z)-r_n(z)=h_nz^{2n+1}+O(z^{2n+2}),
 \qquad \pi_n(x)=x^nq_n(1/x).
\]
Coefficient comparison proves \eqref{ert:pade-orthogonality};
it also gives, for $n\ge\ell\ge0$,
\begin{equation}\label{ert:constant-pade-row}
 \sum_{t=0}^np_{nt}\mu_{t-\ell}=c_n
 \frac1{(2\ell+1)!}
 \prod_{r=0}^{2\ell-1}(2n-r)(2n+3+r).
\end{equation}
Specifically, this sum is $[z^{n-\ell}]r_n(z)$, and its
factorial coefficient formula is the right side. Thus both
$c_n$ and $h_n$ have been obtained without an assumed Pad\'e
normality assertion.

Replace row $n$ of $H+U^{\mathrm{pert}}$ by the linear combination
with coefficients $p_{nt}$, for $0\le n<b$. This simultaneous
transformation is unit lower triangular. In columns $k\ge a$ the
result is $\mathcal N(x^{k-a}\pi_n)$. Hence the final $a$ rows
vanish there, while the first $b-a$ rows in those columns form an
upper triangular square matrix with diagonal
$h_0,\ldots,h_{b-a-1}$. The retained $a$-square matrix is
\[
 A^{\mathrm{ret}}_{ik}=
 \sum_{t=0}^{n_i}p_{n_it}\mu_{t+k-a+1}
 \left(1+\frac{(2t+1)!(2b-2k)!}
 {(2a-2k)!(2b+2t-2a+1)!}\right).
\]
Moving the first $a$ columns past the others contributes
$(-1)^{a(b-a)}$. Reverse the retained columns, putting
$\ell=a-1-k$, and divide row $i$ by $c_{n_i}$.
This adds the determinant factor
$(-1)^{a(a-1)/2}\prod_i c_{n_i}$.

It remains to identify this divided matrix with \eqref{ert:A}.
Its unperturbed entry is $F_{i\ell}$ by
\eqref{ert:constant-pade-row}. Writing $n=n_i$, $N=N_{i\ell}$,
and
\[
 W(j)=(j+2\ell+1)(2n+2\ell+j+1)^{\underline{2i}},\qquad
 C=\frac{R_\ell!}{(2n+1)(2n+2)(2\ell+2)!N!},
\]
factorial cancellation gives the perturbation entry
$2C\sum_{r=0}^{n-\ell}\binom N{2r}B_{2r}W(2r)$.
Since $W(N-y)$ is exactly the polynomial in \eqref{ert:newton},
use the finite identity
\[
 \sum_{j=0}^N\binom Nj B_j(N-j)^{\underline h}
    =N^{\underline h}B_{N-h}(1).
\]
It follows immediately by taking $N^{\underline h}$ outside
the binomial sum and using the defining translation formula for
Bernoulli polynomials. Here $N\ge2i+3$, so all indices $N-h$
are at least two. Only odd $h$ survive. Restoring the omitted
$j=1$ term, with $B_1=-1/2$, gives
\[
 \sum_r\binom N{2r}B_{2r}W(2r)
  =\sum_{k=0}^i d_{i,2k+1}N^{\underline{2k+1}}
           B_{2b-2(a-i+\ell+k)}+\frac N2W(1).
\]
Finally, $CNW(1)=F_{i\ell}$ by factorial cancellation. The divided
retained entry is therefore $2F_{i\ell}$ plus the Bernoulli term
in \eqref{ert:A}. Restoring all row and column factors proves
\eqref{ert:exact-multiplier}.
\end{proof}

For reference, the unperturbed base determinant has the exact
normalization
\begin{equation}\label{ert:base-determinant}
 \det H=(-1)^{\epsilon_{a,b}}
 \left(\prod_{n=0}^{b-a-1}h_n\right)
 \left(\prod_{i=0}^{a-1}c_{n_i}\right)\det F\ne0.
\end{equation}
One can verify the last inequality without using it in the lemma.
Conjugate $H$ by $\operatorname{diag}(1,-1,\ldots)$.
Euler's Bernoulli formula shows that entries with $t+k\ge a$
come from a strictly definite positive moment matrix,
with sign $(-1)^a$. For $a=2h$ the resulting block matrix is
$\left(\begin{smallmatrix}0&B\\B^{\mathsf T}&P\end{smallmatrix}\right)$,
where $P$ is positive definite and $B$ has full row rank $h$.
For $a=2h+1$ it is
$\left(\begin{smallmatrix}ee^{\mathsf T}&B\\B^{\mathsf T}&-P\end{smallmatrix}\right)$,
where $P$ is positive definite, $e=(0,\ldots,0,1)^{\mathsf T}$,
and the first $h$ rows of $B$ have a nonsingular antitriangular
minor. The Schur complements are respectively strictly negative
and strictly positive. The inertia is $(b-h,h)$ in the even case
and $(h+1,b-h-1)$ in the odd case. Together with
\eqref{ert:base-determinant}, this also proves $\det F>0$.
Strict positivity of the moment blocks follows from
$\sum_{m\ge1}(2\pi m)^{-2s}P_0((2\pi m)^{-2})^2>0$ for a
nonzero polynomial $P_0$ and $s\ge1$.

\subsection{The general highest-degree determinant identity}

All denominators can be cleared constructively. Set
\begin{align*}
 J_a(b)&=\prod_{i=0}^{a-1}(2n_i+1)(2n_i+2),\\
 E_a(b)&=\prod_{i,\ell=0}^{a-1}
           \prod_{r=1}^{\max\{0,2i-4\ell-1\}}(R_\ell+r),\\
 D_a(b)&=(2a)!J_a(b)E_a(b),\qquad
 e_a=\sum_{i,\ell=0}^{a-1}\max\{0,2i-4\ell-1\}.
\end{align*}
The symbols $E_a(b)$ and $D_a(b)$ here are auxiliary polynomials;
they are unrelated to the moment coordinates.
Because $N_{i\ell}-R_\ell=2i-4\ell-1$, the polynomial $D_a$
clears every negative factorial ratio in \eqref{ert:A}, every
row denominator, and every constant factorial denominator.
Consequently
\[
 \widetilde{\mathcal A}=D_a\mathcal A
       \in\mathbb Z[b,U_1,\ldots,U_{2a-1}]^{a\times a},
 \qquad \mathscr P_a=\det\widetilde{\mathcal A}
       \in\mathbb Z[b,U_1,\ldots,U_{2a-1}].
\]
Also $D_a(b)>0$ for $b\ge2a$.

Put
\[
 F_0=(2b-4a+2)!,\qquad
 V_j=U_j\frac{F_0}{(2b-2j)!},\qquad
 \kappa_a=\frac{a!((2a)!)^a}{\prod_{\ell=0}^{a-1}(2\ell+2)!},
\]
and
\begin{align}
 \Gamma_a(b)={}&\kappa_aJ_a(b)^{a-1}E_a(b)^a
 \prod_{i=0}^{a-1}(4n_i+2)^{\underline{2i}}
 \prod_{\ell=0}^{a-1}\prod_{r=1}^{2a+2\ell}(2b-4a+2+r).
 \label{ert:Gamma}
\end{align}
The subscript $[a]$ denotes the part homogeneous of total degree
$a$ in the $U$ variables. The exact factorization is
\begin{equation}\label{ert:leading}
 (\mathscr P_a)_{[a]}
 =(-1)^{a(a+1)/2}\Gamma_a(b)
           \det(V_{i+j+1})_{0\le i,j<a}.
\end{equation}
To prove it, write $T=\mathcal A-2F$ and substitute the $V$ variables.
For the term $U_{a-i+\ell+k}$,
$N_{i\ell}-(2k+1)=2b-2(a-i+\ell+k)$, so all nonconstant
factorials cancel. Exactly as in the order-eight proof,
\[
 T=\operatorname{diag}\!\left(\frac1{(2n_i+1)(2n_i+2)}\right)
    L\,(V_{a-i+\ell})_{0\le i,\ell<a}
    \operatorname{diag}\!\left(\frac{R_\ell!}{F_0(2\ell+2)!}\right),
\]
where $L_{i,i-k}=d_{i,2k+1}$ and
$d_{i1}=-(i+1)(4n_i+2)^{\underline{2i}}$.
The $a$ diagonal signs of $L$, together with row reversal of the
middle matrix, give $(-1)^{a(a+1)/2}$. The remaining products are
precisely \eqref{ert:Gamma} after multiplication by $D_a^a$.
This proves \eqref{ert:leading} for every $a$, without evaluating
any large determinant experimentally.

At $U_j=2B_{2b-2j}$, Euler's Bernoulli formula gives
\begin{equation}\label{ert:actual-leading}
 \left|(\mathscr P_a)_{[a]}\right|
 =\frac{\Gamma_a(b)(4F_0)^a}{(2\pi)^{2a(b-a)}}
   \det\bigl(\zeta(2b-4a+2+2i+2j)\bigr)_{0\le i,j<a}.
\end{equation}
Its sign is $(-1)^{ab+a(a+1)/2}$. Define the positive rational
number
\[
 V_a=\prod_{1\le r<s\le a}(r^{-2}-s^{-2})^2,
 \qquad V_1=1.
\]
Keeping the first $a$ support points in Cauchy--Binet proves
\begin{equation}\label{ert:moment-lower}
 \det\bigl(\zeta(2b-4a+2+2i+2j)\bigr)
       \ge V_a(a!)^{-2b+4a-2}>0.
\end{equation}
Thus the highest-degree term is always nonzero in the stated range.

\subsection{An exact algorithm for a sufficient threshold}

Let $\|f\|_1$ denote the sum of the absolute values of all integer
coefficients of a polynomial in $b,U_1,\ldots,U_{2a-1}$. Compute
\[
 E_a^*=\max_{0\le i,\ell<a}\|
                  \widetilde{\mathcal A}_{i\ell}\|_1,
 \qquad K_a=a!(E_a^*)^a,
 \qquad g_a=6a^2-4a+ae_a.
\]
These are finite integer calculations on only $a^2$ explicitly given
entries; no expansion of the whole determinant is required.
The coefficient norm is submultiplicative, so the full coefficient
norm of $\mathscr P_a$ is at most $K_a$.

Each scaled entry has degree at most $6a+e_a-4$ in $b$.
For the nonconstant part this follows from
$\deg_b d_{ih}\le2i+1-h$, the degree $h$ of
$N_{i\ell}^{\underline h}$, and the degree difference
$R_\ell-N_{i\ell}=1+4\ell-2i$; the total bound is
$\deg D_a+4\ell$. The same bound holds for $D_aF_{i\ell}$.
Hence $\deg_b\mathscr P_a\le g_a$.
The factored polynomial \eqref{ert:Gamma} also has degree exactly
$g_a$. Every linear factor in it is at least $b$ when $b\ge4a$,
so
\[
 \Gamma_a(b)\ge\kappa_a b^{g_a}\quad(b\ge4a).
\]

For $b\ge4a+10$, define
\[
 T_{a,b}=\frac{8(2b)^{4a-4}F_0}{(2\pi)^{2b-4a+2}}\ge1.
\]
For the last inequality, put $s_0=2b-4a+2\ge26$;
$26!>8^{26}$ implies $s_0!>8^{s_0}$ by induction,
and $2\pi<8$.
As before, $|2B_{2b-2j}|\le T_{a,b}$ for $1\le j\le2a-1$.
Since the remaining terms have total $U$ degree at most $a-1$,
\[
 |(\mathscr P_a)_{[\le a-1]}|
      \le K_a b^{g_a}T_{a,b}^{a-1}.
\]
Combining this with \eqref{ert:actual-leading} and
\eqref{ert:moment-lower}, cancelling the common powers of $2\pi$,
and then using $1<2\pi<8$, gives
\begin{equation}\label{ert:bound}
 \frac{|(\mathscr P_a)_{[\le a-1]}|}{|(\mathscr P_a)_{[a]}|}
 \le R_a(b):=C_a\frac{b^{L_a}(8a!)^{2b}}{(2b-4a+2)!},
\end{equation}
where every constant is explicit and rational:
\begin{align}
 L_a&=4(a-1)^2,\qquad \chi_a=2a^2-6a+2,\notag\\
 C_a&=\frac{K_a\,8^{a-1+\max(\chi_a,0)}2^{L_a}}
 {\kappa_a4^aV_a(a!)^{4a-2}}.
 \label{ert:constant}
\end{align}
The use of $8^{\max(\chi_a,0)}$ also covers $a=1,2$, where
$\chi_a$ is negative.

Here is a terminating purely arithmetic prescription. Start at
\[
 B_a^{(0)}=\max\{4a+10,\ L_a,\ 16a!\},
\]
and increase the integer $b$ until the exact rational comparison
$R_a(b)<1$ holds. Define $B_a$ to be that first integer. The procedure
must terminate, because for every $b\ge B_a^{(0)}$,
\[
 \frac{R_a(b+1)}{R_a(b)}
 =\left(1+\frac1b\right)^{L_a}
 \frac{(8a!)^2}{(2b-4a+4)(2b-4a+3)}<\frac34.
\]
Indeed $b\ge L_a$ gives $(1+1/b)^{L_a}<3$,
$b\ge4a$ makes both denominator factors at least $b$, and
$b\ge16a!$ supplies the last numerical factor.
The comparison uses integer factorials and powers after clearing
the fixed rational denominator of $C_a$.

For every $b\ge B_a$, \eqref{ert:bound} is strictly less than one,
so the highest-degree term cannot be cancelled. Equations
\eqref{ert:equivalence} and $D_a(b)>0$ prove
Theorem~\ref{ert:main}. In particular, for each input $a$ one may
inspect the finite set $a<b<B_a$ exactly: this either proves the
entire fixed-order family or produces a genuine counterexample.
The construction does not perform those finite calculations for all
$a$, and does not settle the unrestricted conjecture.

\begin{table}[htbp]
\centering
\begin{tabular}{rrrr}
\toprule
$a$ & Smaller order $p=2a$ & Sufficient $B_a$ & Guaranteed even $q$ \\
\midrule
1 & 2  & 16   & $q\ge32$ \\
2 & 4  & 45   & $q\ge90$ \\
3 & 6  & 142  & $q\ge284$ \\
4 & 8  & 467  & $q\ge934$ \\
5 & 10 & 1920 & $q\ge3840$ \\
\bottomrule
\end{tabular}
\caption{Exact sufficient thresholds returned by the constructive
algorithm. No minimality is asserted.}
\label{ert:threshold-table}
\end{table}

The first four rows of Table~\ref{ert:threshold-table} are exact
diagnostics of the general construction; complete-family results
already establish reconstruction for every larger even order in those
cases, including the order-eight theorem proved here. The fifth row
gives a concrete further consequence: for smaller order $p=10$,
optimal-prefix reconstruction holds for every even $q\ge3840$.
The algorithm computes $B_5=1920$ by an exact rational comparison.
It does not check the remaining finite range $6\le b<1920$.
